%=============================================================================!
% 3.13  A LITHIUM ATOM ON A MODEL SURFACE                                     !
%=============================================================================!
\section{A lithium atom on a model surface}
\label{sec:en-surface}

The bead on its wire showed what an energy diagram says in one dimension.
This section draws the same picture for an atom moving in two: a lithium
atom on a layer of carbon atoms arranged in hexagons, as in a sheet of
graphite. We build a simple model of its potential energy, check that
its force does the same work along every path, and follow the atom as it
hops from site to site with its total energy fixed.

\subsection*{The model}

A lithium atom resting on a carbon layer feels a potential energy that
depends on where it sits relative to the carbon atoms. We follow only its
motion parallel to the layer, at a fixed height, with the carbon atoms
held still, so $U$ is a function of the two coordinates $\vb{r} = (x,
y)$. Our model has three kinds of site (\cref{fig:en-surface}a). Above
the centre of a hexagon, a hollow, we make the energy lowest. Above a
carbon atom, a top, we make it highest. Midway between two neighbouring hollows, above the middle
of the bond shared by their hexagons, lies a bridge, the lowest point of
the ridge between the two hollows, and so a saddle point
(\cref{sec:en-gradient}). We take the spacing of neighbouring hollows to
be $a = \SI{2.46}{\angstrom}$, chosen to resemble the hexagons of a layer
of graphite, and the rise from a hollow to a bridge, the barrier, to be $U_{\mathrm b} =
\SI{0.3}{\electronvolt}$. Both are chosen to make the model easy to read;
neither is fitted to any measurement or calculation.

A simple function with this pattern is a sum of three ripples. A single
ripple $\cos(\vb{g}\cdot\vb{r})$, for a fixed vector $\vb{g}$, is
constant along every line perpendicular to $\vb{g}$, since moving along
such a line leaves $\vb{g}\cdot\vb{r}$ unchanged, and it rises and falls
like a sheet of corrugated iron as $\vb{r}$ moves along $\vb{g}$, with
crests $2\pi/\lvert\vb{g}\rvert$ apart. Three such ripples whose vectors
$\vb{g}_1$, $\vb{g}_2$ and $\vb{g}_3$ point at \ang{-30}, \ang{90} and
\ang{210}, all of length $\lvert\vb{g}\rvert = 4\pi/(\sqrt3\,a) =
\SI{2.949}{\per\angstrom}$, so that their crests lie $2\pi/\lvert\vb{g}
\rvert = \sqrt3\,a/2 = \SI{2.130}{\angstrom}$ apart, the distance $a\sin
\ang{60}$ between neighbouring rows of hollows, cross to make a hollow
above the centre of every hexagon, and we set
\begin{equation}
   U(\vb{r}) = \frac{U_{\mathrm b}}{4}\,\Bigl[\,3 - \cos(\vb{g}_1\cdot\vb{r})
   - \cos(\vb{g}_2\cdot\vb{r}) - \cos(\vb{g}_3\cdot\vb{r})\Bigr] .
   \label{eq:en-surface}
\end{equation}
At the origin all three cosines are 1, and $U = 0$: the origin is a
hollow. At the bridge $(a/2, 0)$, using $\vb{g}_k\cdot\vb{r} =
\lvert\vb{g}\rvert\lvert\vb{r}\rvert\cos\theta_k$ with $\theta_k$ the
angle between them (\cref{sec:mo-plane}), and $\lvert\vb{g}\rvert\,a/2 =
2\pi/\sqrt3$,
\begin{align*}
   \vb{g}_1\cdot\vb{r} &= \tfrac{2\pi}{\sqrt3}\cos\ang{-30} = \pi , &
   \vb{g}_2\cdot\vb{r} &= \tfrac{2\pi}{\sqrt3}\cos\ang{90} = 0 , &
   \vb{g}_3\cdot\vb{r} &= \tfrac{2\pi}{\sqrt3}\cos\ang{210} = -\pi ,
\end{align*}
since $\cos\ang{-30} = \sqrt3/2$ and $\cos\ang{210} = -\sqrt3/2$. The
cosines are $-1$, $1$ and $-1$, and $U = (U_{\mathrm b}/4)\bigl(3 - (-1) - 1 -
(-1)\bigr) = (U_{\mathrm b}/4)\times4 = U_{\mathrm b}$. The top $(a/2, a/(2\sqrt3))$
lies a distance $\sqrt{a^2/4 + a^2/12} = \sqrt{a^2/3} = a/\sqrt3 =
\SI{1.420}{\angstrom}$ from the hollow, in the direction at \ang{30} to the
$x$ axis, since $\tan\theta = \bigl(a/(2\sqrt3)\bigr)/(a/2) = 1/\sqrt3$.
There $\lvert\vb{g}\rvert\,a/\sqrt3 = 4\pi/3$, and the angles from
$\vb{r}$ to $\vb{g}_1$, $\vb{g}_2$ and $\vb{g}_3$ are \ang{60}, \ang{60}
and \ang{180}, so the three products are $\frac{4\pi}{3}\cos\ang{60} =
2\pi/3$, $2\pi/3$ and $\frac{4\pi}{3}\cos\ang{180} = -4\pi/3$. Each cosine
is $-\frac12$ (\cref{ex:en-top} checks this in components), and $U = (U_{\mathrm b}/4)(3 + \frac32) = 9U_{\mathrm b}/8 =
\SI{0.3375}{\electronvolt}$. Shifting $\vb{r}$ by $a$ along $x$ changes
$\vb{g}_1\cdot\vb{r}$ by $2\pi$, $\vb{g}_2\cdot\vb{r}$ by 0 and
$\vb{g}_3\cdot\vb{r}$ by $-2\pi$, which leaves every cosine unchanged, and
the same holds for a shift by $a$ at \ang{60}, so the pattern repeats from
hollow to hollow.

The force follows from \cref{eq:en-gradient}. Writing $\vb{g}\cdot\vb{r} =
g_x x + g_y y$, the chain rule gives $\partial\cos(\vb{g}\cdot\vb{r})/
\partial x = -g_x\sin(\vb{g}\cdot\vb{r})$, and likewise with $g_y$ for
$y$, so $\grad\cos(\vb{g}\cdot\vb{r}) = -\sin(\vb{g}\cdot\vb{r})\,\vb{g}$,
and
\begin{align*}
   \grad U &= \frac{U_{\mathrm b}}{4}\sum_{k=1}^{3}\sin(\vb{g}_k\cdot\vb{r})\,
      \vb{g}_k
      && \text{the 3 is constant, and each $-\cos$ gives $+\sin$,}\\
   \vb{F} = -\grad U &= -\frac{U_{\mathrm b}}{4}\sum_{k=1}^{3}
      \sin(\vb{g}_k\cdot\vb{r})\,\vb{g}_k .
\end{align*}
At the hollow and the bridge every sine is zero, since $\sin 0 =
\sin(\pm\pi) = 0$. At the top the three sines are equal, because
$\sin(-4\pi/3) = \sin(-4\pi/3 + 2\pi) = \sin(2\pi/3)$, so the force is
proportional to $\vb{g}_1 + \vb{g}_2 + \vb{g}_3$, which is zero: the $x$
components of the three directions add to $\cos\ang{-30} + \cos\ang{90} +
\cos\ang{210} = \sqrt3/2 + 0 - \sqrt3/2 = 0$, and the $y$ components to
$-\frac12 + 1 - \frac12 = 0$. Every site is an equilibrium. The function \texttt{energy.hexagonal\_surface} of
\texttt{mdlab} returns $U$ and $\vb{F}$ together, and its tests check the
force against central differences of $U$ and the site energies above.

\subsection*{Work on the surface}

\Cref{fig:en-surface}(a) shows the contours of $U$, the force, and two
paths from the hollow $A$ at the origin to the top $B$ at $(3.69,
0.710)$\,\si{\angstrom}, one spacing $a$ along $x$ from the top nearest
$A$: a straight line, and an arc that bulges to its left, lying
$1.2\sin(\pi s)$\,\si{\angstrom} from the line at the point a fraction $s$
of the way along it, so \SI{1.2}{\angstrom} at the middle. The function \texttt{work} gives
\SI{-0.3375}{\electronvolt} along both, equal to $U(A) - U(B)$, as a
conservative force requires. In \cref{fig:en-surface}(b) the swirl of
\cref{eq:en-swirl}, centred on $A$ with $c = \SI{0.005}{\electronvolt
\per\angstrom\squared}$, is added to the force. Along the straight path,
which runs directly away from the centre of the swirl, the swirl is
perpendicular to the path and does no work, so the work stays at
\SI{-0.3375}{\electronvolt}. Along the arc it is
\SI{-0.3662}{\electronvolt}. The arc followed by the straight path
backwards is a clockwise loop, and the curl of the total force is $0 +
2c$, the surface force being a gradient, so by Green's theorem the work
round this loop, $W_{\mathrm{arc}} - W_{\mathrm{straight}}$, is $-2c$
times the area between the paths. That area is the bulge integrated along
the chord, of length $\lvert AB\rvert = \sqrt{3.69^2 + 0.710^2} =
\SI{3.758}{\angstrom}$, on which a step $\dd s$ in the fraction moves
$\lvert AB\rvert\,\dd s$,
$\int_0^1 1.2\sin(\pi s)\,\lvert AB\rvert\,\dd s = 2\times1.2\times
3.758/\pi = \SI{2.871}{\angstrom\squared}$, using $\int_0^1\sin(\pi s)\,
\dd s = [-\cos(\pi s)/\pi]_0^1 = 2/\pi$; and $2c\times2.871 =
\SI{0.0287}{\electronvolt}$ is exactly the difference between the two
works.

\begin{figure}[htb]
\centering
\includegraphics{ch03_energy/surface}
\caption{Work on the model surface. Open circles mark hollows and dark
dots the carbon atoms, at the tops. (a) Contours of $U$ at 0.05, 0.10,
0.15, 0.20, 0.25 and \SI{0.32}{\electronvolt}, the force $-\grad U$
(grey arrows), and two paths from the hollow $A$ to the top $B$; the work
is \SI{-0.3375}{\electronvolt} along both. (b) The same paths after a
swirl with $c = \SI{0.005}{\electronvolt\per\angstrom\squared}$, centred
on $A$, is added to the force (grey arrows). No potential energy exists to
contour, and the works differ by \SI{0.0287}{\electronvolt}.}
\label{fig:en-surface}
\end{figure}

\subsection*{Hopping with a fixed energy}

A lithium atom leaves the hollow at the origin with \SI{0.33}{
\electronvolt} of kinetic energy, \SI{10}{\percent} more than the barrier
$U_{\mathrm b}$, moving at \ang{37.5} to the $x$ axis. By \cref{eq:en-mv2} its
speed is
\begin{equation*}
   v = \sqrt{\frac{2K}{m\times103.6427}}
   = \sqrt{\frac{2\times0.33}{6.94\times103.6427}}
   = \SI{0.03029}{\angstrom\per\femto\second} ,
\end{equation*}
about \SI{3000}{\metre\per\second}. Its total energy, $E =
\SI{0.33}{\electronvolt}$, exceeds $U$ everywhere except in small islands
around each carbon atom, where $U$ rises to \SI{0.3375}{\electronvolt}:
the atom may go anywhere on the surface except into these islands, just
as the bead at $h_E = \SI{1.3}{\metre}$ could go anywhere except up the ends
of its track. \Cref{fig:en-li} shows the motion over
\SI{1.5}{\pico\second}, or \SI{1500}{\femto\second} (a picosecond is
\SI{1000}{\femto\second}), computed by a computer solution of the equations
of motion, a method from SciPy, a library of scientific routines for Python, that
advances the state in
short steps, as in \cref{sec:ne-equations}, but combines several force
evaluations per step and adjusts each step so that the error it
estimates for that step stays below about one part in $10^{11}$; over
many steps these small errors add up. Chapter~9 develops the methods that
molecular dynamics itself uses. The atom crosses from hollow to hollow
four times, visiting five hollows, while its kinetic energy rises and falls between
\num{0.0005} and \SI{0.33}{\electronvolt}. Its total energy stays at
\SI{0.33}{\electronvolt} to within \SI{2.5e-11}{\electronvolt}, which is
the error of the computer solution.

\begin{figure}[htb]
\centering
\includegraphics{ch03_energy/li_energy}
\caption{A lithium atom of \SI{6.94}{\amu} on the model surface, leaving
the hollow at the origin with \SI{0.33}{\electronvolt} of kinetic energy.
(a) The potential energy, shaded from the hollows (light) to the bridges
and tops (dark), the boundary $U = E$ of the islands around the carbon
atoms (dots), where the atom may not go, and its path over
\SI{1.5}{\pico\second}. (b) Kinetic, potential and total energy over the
first \SI{600}{\femto\second}.}
\label{fig:en-li}
\end{figure}

With less energy the picture changes as it did on the track. Below the
barrier, $E < \SI{0.3}{\electronvolt}$, the allowed region breaks into
separate patches, one around each hollow, and an atom stays in the hollow
it starts in, rattling about its bottom. To hop to a neighbouring hollow
it needs at least the barrier as kinetic energy at the bottom of the
hollow, a speed there of at least $\sqrt{2\times0.3/(6.94\times103.6427)}
= \SI{0.02888}{\angstrom\per\femto\second}$. Real atoms exchange energy
with their neighbours all the time, and Chapter~12 gives the odds that one
collects as much as the barrier; Chapters~16 and 17 measure how fast lithium spreads through graphite by
such hops.

\notebook{03}{launch the lithium atom with your own energy and direction}

\begin{selfcheck}
A lithium atom at the bottom of a hollow of the model surface has
\SI{0.25}{\electronvolt} of kinetic energy. Can it reach a neighbouring
hollow? How fast is it moving where $U = \SI{0.1}{\electronvolt}$?
\end{selfcheck}
